\section{Arms and Controls}
\label{sec:method}

%TL;DR Every arm shares one builder, one build language, one renderer and one canvas, and differs from its control in exactly one ingredient.
All arms share one builder model (\gptmini{}), one build language, one textured renderer and one
$32^3$ canvas, and each arm is a recipe file that differs from its control in one ingredient.
\Cref{tab:arms} lists the ingredients. This section describes the build language
(\cref{sec:language}), the loop without an image (\cref{sec:classic}), the imagination loop
(\cref{sec:dream}), the caption (\cref{sec:caption}), the controls (\cref{sec:controls}) and the
measures that keep the comparison fair (\cref{sec:fair}).

\begin{table}[t]
  \centering
  \small
  \setlength{\tabcolsep}{3pt}
  \begin{tabular}{@{}lcccr@{}}
    \toprule
    Arm & Image & Target text & Turns & \$/build \\
    \midrule
    \classicm{}  & no  & none                        & $1{+}3$        & 0.045 \\
    \dream{}     & yes & none                        & $1{+}{\le}3$   & 0.040 \\
    \dreamcap{}  & yes & own caption & $1{+}{\le}3$   & 0.048 \\
    \midrule
    \shufcap{}   & yes & other caption & $1{+}{\le}3$   & \tbd{} \\
    \textspec{}  & no  & spec & $1{+}3$        & \tbd{} \\
    \bestoffour{}& no  & none                        & $4{\times}1$   & \tbd{} \\
    \bottomrule
  \end{tabular}
  % $/build: card run_20261109_121625 cost table (cost_usd, n=133); T51 reports the same 0.045 / 0.040 / 0.048
  \caption{\textbf{Arms.} ``Turns'' counts one repair turn for execution errors plus the critique
  or match turns; \bestoffour{} draws four independent builds with one repair turn each. All arms
  select the best attempt they produce. The lower three rows are controls; their cost is
  \TODO{n=256 card}.}
  \label{tab:arms}
\end{table}

\subsection{Build language}
\label{sec:language}

%TL;DR The builder writes a short program in a WorldEdit-style language whose commands place many blocks each, over a fixed palette of about seventy block names.
The builder writes a program of one command per line in a language modelled on the in-game
editor WorldEdit~\citep{worldedit}. Commands take inclusive integer coordinates with $y$ up and
the origin at a canvas corner: \texttt{set} places one block, \texttt{fill}, \texttt{clear},
\texttt{box}, \texttt{walls} and \texttt{replace} act on cuboids, \texttt{line},
\texttt{sphere}, \texttt{cylinder}, \texttt{pyramid} and \texttt{gable} draw primitives, and
\texttt{copy}, \texttt{paste} and \texttt{stack} repeat parts. Blocks are named from a curated
palette of about seventy modern names, with block states such as stair orientation where the game
stores them. The system prompt lists the commands and the palette, gives one example program, and
states player scale: a player is two blocks tall and one block wide, so doors are two blocks tall.

%TL;DR Parsing and execution report errors per line instead of failing, which gives every repair turn a precise list of what to fix.
The executor parses the whole program before running it and reports syntax errors with line
numbers. A command that fails is recorded and skipped while the rest executes, writes outside the
canvas are clipped and counted, an unknown block name is an error on its line rather than a
silent substitution, and a command that changes no block raises a warning, which almost always
marks a coordinate mistake. Every arm spends one repair turn on these reports after its first
program.

\subsection{The loop without an image: \classicm{}}
\label{sec:classic}

%TL;DR The compute-matched control writes a program, repairs execution errors, then critiques and rewrites from two renders of its own build three times, and keeps the best attempt.
\classicm{} is the control for every arm that draws an image. After the first program and its
repair turn, it runs three critique turns. In each, the builder sees renders of its current build
from two fixed cameras and rewrites the whole program. A new attempt is accepted when its score
improves and the share of blocks in its largest connected component drops by at most 0.1, which
blocks the cheap trick of adding floating decoration. The arm returns its best attempt under a
score that does not depend on any image target: the similarity between the prompt and the render
under CLIP ViT-B/16~\citep{radford2021clip}, plus half the largest-component share.

\subsection{The imagination loop: \dream{}}
\label{sec:dream}

%TL;DR The dream loop draws one isometric target image, builds toward it, and then compares render and target in up to three match turns.
\dream{} adds a target image, the dream, and nothing else (\cref{alg:dream}). The image model
(\imgmini{}~\citep{openai2025gptimage}, low quality, $1024{\times}1024$) draws the prompt from a
fixed isometric camera in a pinned style: a single build of full cube blocks in flat daylight on a
white background, with no terrain, sky, people or text. The template also states the canvas size
and player scale. The builder sees the dream labelled as the target in its first turn. After the
repair turn, up to three match turns follow. In each, the builder sees the target next to a render
of its build from the dream's camera, labelled as its own, and rewrites the program.

\begin{algorithm}[t]
  \caption{\dream{}, \dreamcap{} and \shufcap{}. Lines marked $\dagger$ run only in the caption arms.}
  \label{alg:dream}
  \small
  \begin{algorithmic}[1]
    \Require prompt $p$, canvas $32^3$, at most $K{=}3$ match turns
    \State $I \gets \textsc{ImageModel}(\text{style}, p)$ \Comment{cached by request}
    \State $c \gets \textsc{Captioner}(I)$, or another prompt's caption \Comment{$\dagger$}
    \State $P_0 \gets \textsc{Builder}(p, I, [c])$; repair once from executor errors
    \For{$k = 1, \dots, K$}
      \State \textbf{if} $s(P_{k-1}) \ge 0.85$ or two turns without gain \textbf{then break}
      \State $R \gets$ render of $P_{k-1}$ at the camera that best matches $I$
      \State $P_k \gets \textsc{Builder}(\text{target } I, \text{yours } R, [c])$
      \State accept $P_k$ if $s$ improves and connectivity holds
    \EndFor
    \State \Return $\arg\max_k s(P_k)$
  \end{algorithmic}
\end{algorithm}

%TL;DR The in-loop score compares render and dream with backbones disjoint from the evaluation's, gated on connectivity, and the loop stops early on success or stall.
The loop scores each attempt against the dream with the cosine similarity of SigLIP-B
embeddings~\citep{zhai2023siglip} (weight 1.0), a silhouette overlap (weight 0.5) and the
largest-component share (weight 0.3); an attempt whose largest component holds less than half its
blocks scores zero. The render is taken from whichever of the four isometric cameras matches the
dream best, because the image model's camera is approximate. The loop stops when the score reaches
0.85 or after two turns without improvement, applies the same acceptance rule as \classicm{}, and
returns the best attempt. Neither in-loop backbone (SigLIP-B here, CLIP ViT-B/16 in \classicm{})
appears among the evaluation's backbones~\citep{anonymous2027blindness}, so no arm optimizes the
quantity it is ranked by.

\subsection{Describing the dream: \dreamcap{}}
\label{sec:caption}

%TL;DR The caption arm adds one vision turn in which a captioner, who never sees the prompt, describes the dream in a fixed schema that the builder reads in every turn.
\dreamcap{} adds one vision turn to \dream{} and changes nothing else. A captioner (also
\gptmini{}) receives the dream and the canvas size, but not the prompt, and describes what the
image shows in a fixed JSON schema: a one-sentence summary, the size and footprint in blocks, the
storey count, the roof shape and material, the palette blocks of the walls, trim, roof, windows,
door and foundation, the openings, distinctive features, notes on the front and side faces,
anything around the building to leave out, and a list of what is uncertain. Fields the image does
not show, such as the back or the interior, are marked unknown. The builder reads the caption as
a description of the target in its first turn and in every match turn. Because the image cache
returns the same dream for the same request, \dreamcap{} and \dream{} build toward identical
images, and any difference between them is the caption.

\subsection{Controls}
\label{sec:controls}

%TL;DR Three controls ask whether the caption helps through its content, whether the image is needed for that content, and whether more samples alone do as well.
Each control answers one question about the caption result.
\begin{itemize}[leftmargin=*,itemsep=1pt,topsep=2pt]
  \item \shufcap{} makes the same calls as \dreamcap{} but shows the builder the caption of the
    dream of the previous prompt in the same arm; the first prompt receives the one of five fixed
    library captions that shares the fewest words with it. Since the captioner never sees the
    prompt, a shown caption carries no trace of where it came from. If \shufcap{} matches
    \dreamcap{}, the gain is extra text, not description.
  \item \textspec{} writes a specification of the house from the prompt alone, in the fields
    of the caption without the size and per-face notes, plus the floor material and storey heights, and runs the \classicm{}
    loop with the specification in place of an image. If \textspec{} matches \dreamcap{}, the
    image is not needed: a structured description is. \TODO{\textspec{} has no recipe yet in
    \texttt{configs/agentic/recipes}; the closest is \texttt{dream-describe-before}, which draws
    an image from the spec. Add a no-image recipe and confirm this description matches it.}
  \item \bestoffour{} draws four independent \classicm{} builds with no critique turns and keeps
    the best under the \classicm{} score. If it matches \dream{} or \dreamcap{}, the gain is
    sampling, not imagination.
\end{itemize}

\subsection{Fair comparison}
\label{sec:fair}

%TL;DR Shared caches, recipe hashes and a cost ledger make arms paired on images, exactly reproducible, and comparable in dollars.
Three mechanisms keep the arms comparable. First, dreams and model replies sit in
content-addressed caches, so arms that issue the same request receive the same response, and a
rerun in read-only mode reproduces a whole arm at zero cost. Second, each recipe resolves to a
hash over all of its fields, and two runs count as the same arm only when their hashes match.
Third, a ledger prices every call before it is made, enforces a budget per build and per run, and
records both the price paid and the full price of calls served from the cache.
All arms run on the same prompts, the titles of held-out houses written as
``\textit{title} (a Minecraft \textit{category})'', so each build has a real house with the same
name in the evaluation set (\cref{sec:setup}).
\TODO{n=256: confirm the 123 added prompts are validation-split titles in the same format, and that the card reports full-price cost.}
